

Par
ailleurs, @@M0@@ définit une section @@M1@@ de @@M2@@:
@@M3@@
Soient @@M4@@ les idéaux définissant @@M5@@ ; on a : @@M6@@, et l'immersion @@M7@@ est régulière; de plus, @@M8@@ est définie par l'idéal @@M9@@ (car @@M10@@ est plat), et son faisceau conormal est donc @@M11@@. On obtient alors la suite exacte :
@@M12@@
Nous allons montrer :

a) que le diagramme (2.2.1) est commutatif ; on obtiendra donc ainsi un morphisme de complexes @@M13@@ du complexe @@M14@@ défini par @@M15@@ dans le complexe @@M16@@ défini par @@M17@@ ;

b) que l'homomorphisme : @@M18@@ défini par (2.2.2), (2.2.3) et (2.2.1) est un isomorphisme ; cela entraînera que le mapping cylinder de @@M19@@ est acyclique, donc que @@M20@@ est un quasi-isomorphisme, donc définit un isomorphisme dans la catégorie dérivée.

Les propriétés a) et b) sont de nature locale sur @@M21@@. On peut donc supposer @@M22@@ et @@M23@@ affines. Posons : @@M24@@, @@M25@@, @@M26@@, @@M27@@, et gardons les mêmes lettres pour noter les homomorphismes et les idéaux correspondants. On a le diagramme :
